Cas du mode XOn considere la diffusion quasilineaire a l'harmonique
$n$ par l'onde EC avec une polarisation de mode X de sorte que la
composante predominante du champ electrique soit $E_{\mathbf{k},-}$.
Pour les electrons, \[
\Omega=\frac{-\omega_{ce}}{\gamma}<0\]
 avec \[
\omega_{ce}=\frac{q_{e}B}{m_{e}}\]
 il faut en realite prendre un harmonique $n<0$. Dans ce paragraphe,
$n$ correspond en fait a sa valeur absolue. En notant que $\left|J_{-n}\right|=\left|J_{n}\right|$
, on a alors \[
D_{\mathbf{k}}^{\text{EC-X}(n)}\simeq q_{e}^{2}\left|E_{\mathbf{k},-}\right|^{2}J_{n-1}^{2}\]


On se limite aux cas \[
\frac{k_{\bot}v_{\bot}}{\Omega}\simeq n\frac{v_{\bot}}{c}<1\]
 de sorte que \[
J_{n}(z)\simeq\frac{(z/2)^{n}}{n!}\]
 et \begin{eqnarray*}
D_{\mathbf{k}}^{\text{EC-X}(n)} & \simeq & \frac{1}{(n-1)!^{2}}\left(\frac{k_{\bot}v_{\bot}}{2\Omega}\right)^{2(n-1)}q_{e}^{2}\left|E_{\mathbf{k},-}\right|^{2}\\
 & = & \frac{1}{(n-1)!^{2}}\left(N_{\bot}\frac{\gamma v_{\bot}}{c}\frac{\omega}{2\omega_{ce}}\right)^{2(n-1)}q_{e}^{2}\left|E_{\mathbf{k},-}\right|^{2}\end{eqnarray*}


Ainsi, le coefficient de diffusion quasilineaire pour l'onde EC en
mode X devient \begin{eqnarray*}
D^{\text{EC-X}(n)}(p,\xi) & = & \left(\frac{p_{\bot}}{p_{Te}}\right)^{2(n-1)}\lim_{V\rightarrow\infty}\frac{\pi q_{e}^{2}}{V}\int\frac{d^{3}k}{(2\pi)^{3}}\frac{1}{(n-1)!^{2}}\\
 &  & \times\left(N_{\bot}\frac{p_{Te}}{m_{e}c}\frac{\omega}{2\omega_{ce}}\right)^{2(n-1)}\left|E_{\mathbf{k},-}\right|^{2}\delta\left(\omega-k_{\Vert}v_{\Vert}-\frac{n\omega_{ce}}{\gamma}\right)\\
 & = & \left(\frac{p_{\bot}}{p_{Te}}\right)^{2(n-1)}\frac{v_{Te}}{\left|v_{\Vert}\right|}\lim_{V\rightarrow\infty}\frac{\pi q_{e}^{2}}{V}\int\frac{d^{3}k}{(2\pi)^{3}}\frac{1}{(n-1)!^{2}}\\
 &  & \times\left(N_{\bot}\frac{p_{Te}}{m_{e}c}\frac{\omega}{2\omega_{ce}}\right)^{2(n-1)}\left|E_{\mathbf{k},-}\right|^{2}\delta\left(k_{\Vert}-\frac{\omega}{v_{\Vert}}+\frac{n\omega_{ce}}{\gamma v_{\Vert}}\right)\frac{1}{v_{Te}}\end{eqnarray*}


On considere un rayon EC a une position polo\"{\i}dale particuliere
$\theta=\theta_{b}$, avec un spectre du champ electrique gaussien
en $k_{\Vert}$ centre autour de $k_{\Vert,0}$ avec une demi-largeur
$\Delta k_{\Vert}$ a valeur $1/e$. Dans ces conditions, on peut
modeliser le champ electrique par \[
\left|E_{\mathbf{k},-}\right|^{2}(k_{\Vert},\overline{k}_{\bot})=\left|E_{\mathbf{k},-,0}\right|^{2}(\overline{k}_{\bot})\exp\left(-\frac{\left(k_{\Vert}-k_{\Vert,0}\right)^{2}}{\Delta k_{\Vert}^{2}}\right)2\pi\Delta\theta\delta(\theta-\theta_{b})\]
 de sorte que \[
D^{\text{EC-X}(n)}(p,\xi)=\left(\frac{p_{\bot}}{p_{Te}}\right)^{2(n-1)}\frac{v_{Te}}{\left|v_{\Vert}\right|}\exp\left(-\frac{\left(k_{\Vert,\text{res}}-k_{\Vert,0}\right)^{2}}{\Delta k_{\Vert}^{2}}\right)2\pi\delta(\theta-\theta_{b})D_{0}^{\text{EC-X}(n)}\]
 avec \[
D_{0}^{\text{EC-X}(n)}=\lim_{V\rightarrow\infty}\frac{\pi q_{e}^{2}\Delta\theta}{V}\int\frac{d^{2}k_{\bot}}{(2\pi)^{3}}\frac{1}{(n-1)!^{2}}\left(N_{\bot}\frac{p_{Te}}{m_{e}c}\frac{\omega}{2\omega_{ce}}\right)^{2(n-1)}\left|E_{\mathbf{k},-,0}\right|^{2}(\overline{k}_{\bot})\frac{1}{v_{Te}}\]
 et \[
k_{\Vert,\text{res}}=\frac{\omega}{v_{\Vert}}\left(1-\frac{n\omega_{ce}}{\gamma\omega}\right)\]


Le coefficient de diffusion quasilineaire pour l'onde EC en mode X
devient alors \begin{eqnarray*}
D_{(0)}^{\text{EC-X}(n)}(p,\xi_{0}) & \equiv & \left\{ \Psi D^{\text{EC-X}(n)}(p,\xi)\right\} \\
 & = & \left\{ \Psi\left(\frac{p_{\bot}}{p_{Te}}\right)^{2(n-1)}\frac{v_{Te}}{\left|v_{\Vert}\right|}\exp\left(-\frac{\left(k_{\Vert,\text{res}}-k_{\Vert,0}\right)^{2}}{\Delta k_{\Vert}^{2}}\right)2\pi\delta(\theta-\theta_{b})D_{0}^{\text{EC-X}(n)}\right\} \\
 & = & \left(\frac{p}{p_{Te}}\right)^{2(n-1)}\left(1-\xi_{0}^{2}\right)^{(n-1)}\frac{v_{Te}}{v|\xi_{0}|}\\
 &  & \times\left\{ \frac{\xi_{0}}{\xi}\Psi^{n}\exp\left(-\frac{\left(k_{\Vert,\text{res}}-k_{\Vert,0}\right)^{2}}{\Delta k_{\Vert}^{2}}\right)2\pi\delta(\theta-\theta_{b})D_{0}^{\text{EC-X}(n)}\right\} \\
 & = & \left(\frac{p}{p_{Te}}\right)^{2(n-1)}\left(1-\xi_{0}^{2}\right)^{(n-1)}\frac{v_{Te}}{v|\xi_{0}|}\\
 &  & \frac{1}{\lambda}\left[\frac{1}{2}\sum_{\sigma=\pm1}\right]_{T}\int_{-\theta_{c}}^{\theta_{c}}\frac{d\theta}{2\pi}\frac{\xi_{0}^{2}}{\xi^{2}}\Psi^{n}\exp\left(-\frac{\left(k_{\Vert,\text{res}}-k_{\Vert,0}\right)^{2}}{\Delta k_{\Vert}^{2}}\right)2\pi\delta(\theta-\theta_{b})D_{0}^{\text{EC-X}(n)}\end{eqnarray*}
 et finalement \begin{eqnarray*}
D_{(0)}^{\text{EC-X}(n)}(p,\xi_{0}) & \equiv & \left\{ \Psi D^{\text{EC-X}(n)}(p,\xi)\right\} \\
 & = & \left(\frac{p}{p_{Te}}\right)^{2(n-1)}\left(1-\xi_{0}^{2}\right)^{(n-1)}\frac{v_{Te}}{v|\xi_{0}|}\\
\times\frac{1}{\lambda}\left[\frac{1}{2}\sum_{\sigma=\pm1}\right]_{T}H(|\theta_{b}| & < & \theta_{c})\frac{\xi_{0}^{2}}{\xi_{b}^{2}}\Psi_{b}^{n}\exp\left(-\frac{\left(k_{\Vert,\text{res},b}-k_{\Vert,0}\right)^{2}}{\Delta k_{\Vert}^{2}}\right)D_{0,b}^{\text{EC-X}(n)}\end{eqnarray*}
 avec \[
D_{0,b}^{\text{EC-X}(n)}=\lim_{V\rightarrow\infty}\frac{\pi q_{e}^{2}\Delta\theta_{b}}{V}\int\frac{d^{2}k_{\bot}}{(2\pi)^{3}}\frac{1}{(n-1)!^{2}}\left(N_{\bot,b}\frac{p_{Te}}{m_{e}c}\frac{\omega_{b}}{2\omega_{ce}}\right)^{2(n-1)}\left|E_{\mathbf{k},-,0,b}\right|^{2}(\overline{k}_{\bot})\frac{1}{v_{Te}}\]


Cas du mode O

On considere la diffusion quasilineaire a l'harmonique $n$ par l'onde
EC avec une polarisation de mode O de sorte que la composante predominante
du champ electrique soit $E_{\mathbf{k},\Vert}$. On a alors \[
D_{\mathbf{k}}^{\text{EC-O}(n)}\simeq q_{e}^{2}\frac{p_{\Vert}^{2}}{p_{\bot}^{2}}\left|E_{\mathbf{k},\Vert}\right|^{2}J_{n}^{2}\]


On se limite egalement aux cas \[
\frac{k_{\bot}v_{\bot}}{\Omega}\simeq n\frac{v_{\bot}}{c}<1\]
 de sorte que \[
J_{n}(z)\simeq\frac{(z/2)^{n}}{n!}\]
 et \begin{eqnarray*}
D_{\mathbf{k}}^{\text{EC-O}(n)} & \simeq & \frac{1}{n!^{2}}\left(\frac{k_{\bot}v_{\bot}}{2\Omega}\right)^{2n}q_{e}^{2}\frac{p_{\Vert}^{2}}{p_{\bot}^{2}}\left|E_{\mathbf{k},\Vert}\right|^{2}\\
 & = & \frac{1}{n!^{2}}\left(N_{\bot}\frac{\gamma v_{\bot}}{c}\frac{\omega}{2\omega_{ce}}\right)^{2n}q_{e}^{2}\frac{p_{\Vert}^{2}}{p_{\bot}^{2}}\left|E_{\mathbf{k},\Vert}\right|^{2}\end{eqnarray*}


Ainsi, le coefficient de diffusion quasilineaire pour l'onde EC en
mode O devient \begin{eqnarray*}
D^{\text{EC-O}(n)}(p,\xi) & = & \left(\frac{p_{\bot}}{p_{Te}}\right)^{2n}\frac{p_{\Vert}^{2}}{p_{\bot}^{2}}\\
 &  & \times\lim_{V\rightarrow\infty}\frac{\pi q_{e}^{2}}{V}\int\frac{d^{3}k}{(2\pi)^{3}}\frac{1}{n!^{2}}\left(N_{\bot}\frac{p_{Te}}{m_{e}c}\frac{\omega}{2\omega_{ce}}\right)^{2n}\left|E_{\mathbf{k},\Vert}\right|^{2}\delta\left(\omega-k_{\Vert}v_{\Vert}-\frac{n\omega_{ce}}{\gamma}\right)\\
 & = & \left(\frac{p_{\bot}}{p_{Te}}\right)^{2n}\frac{p_{\Vert}^{2}}{p_{\bot}^{2}}\frac{v_{Te}}{\left|v_{\Vert}\right|}\\
 &  & \times\lim_{V\rightarrow\infty}\frac{\pi q_{e}^{2}}{V}\int\frac{d^{3}k}{(2\pi)^{3}}\frac{1}{n!^{2}}\left(N_{\bot}\frac{p_{Te}}{m_{e}c}\frac{\omega}{2\omega_{ce}}\right)^{2n}\left|E_{\mathbf{k},\Vert}\right|^{2}\delta\left(k_{\Vert}-\frac{\omega}{v_{\Vert}}+\frac{n\omega_{ce}}{\gamma v_{\Vert}}\right)\frac{1}{v_{Te}}\end{eqnarray*}


On considere un rayon EC a une position polo\"{\i}dale particuliere
$\theta=\theta_{b}$, avec un spectre du champ electrique gaussien
en $k_{\Vert}$ centre autour de $k_{\Vert,0}$ avec une demi-largeur
$\Delta k_{\Vert}$ a la valeur $1/e$. Dans ces conditions, on peut
modeliser le champ electrique par \[
\left|E_{\mathbf{k},\Vert}\right|^{2}(k_{\Vert},\overline{k}_{\bot})=\left|E_{\mathbf{k},\Vert,0}\right|^{2}(\overline{k}_{\bot})\exp\left(-\frac{\left(k_{\Vert}-k_{\Vert,0}\right)^{2}}{\Delta k_{\Vert}^{2}}\right)2\pi\Delta\theta\delta(\theta-\theta_{b})\]
 de sorte que \[
D^{\text{EC-O}(n)}(p,\xi)=\left(\frac{p_{\bot}}{p_{Te}}\right)^{2n}\frac{p_{\Vert}^{2}}{p_{\bot}^{2}}\frac{v_{Te}}{\left|v_{\Vert}\right|}\exp\left(-\frac{\left(k_{\Vert,\text{res}}-k_{\Vert,0}\right)^{2}}{\Delta k_{\Vert}^{2}}\right)2\pi\delta(\theta-\theta_{b})D_{0}^{\text{EC-X}(n)}\]
 avec \[
D_{0}^{\text{EC-O}(n)}=\lim_{V\rightarrow\infty}\frac{\pi q_{e}^{2}\Delta\theta}{V}\int\frac{d^{2}k_{\bot}}{(2\pi)^{3}}\frac{1}{n!^{2}}\left(N_{\bot}\frac{p_{Te}}{m_{e}c}\frac{\omega}{2\omega_{ce}}\right)^{2n}\left|E_{\mathbf{k},\Vert,0}\right|^{2}(\overline{k}_{\bot})\frac{1}{v_{Te}}\]
 et \[
k_{\Vert,\text{res}}=\frac{\omega}{v_{\Vert}}\left(1-\frac{n\omega_{ce}}{\gamma\omega}\right)\]


Le coefficient de diffusion quasilineaire pour l'onde EC en mode O
devient alors \begin{eqnarray*}
D_{(0)}^{\text{EC-O}(n)}(p,\xi_{0}) & \equiv & \left\{ \Psi D^{\text{EC-O}(n)}(p,\xi)\right\} \\
 & = & \left\{ \Psi\left(\frac{p_{\bot}}{p_{Te}}\right)^{2n}\frac{p_{\Vert}^{2}}{p_{\bot}^{2}}\frac{v_{Te}}{\left|v_{\Vert}\right|}\exp\left(-\frac{\left(k_{\Vert,\text{res}}-k_{\Vert,0}\right)^{2}}{\Delta k_{\Vert}^{2}}\right)2\pi\delta(\theta-\theta_{b})D_{0}^{\text{EC-O}(n)}\right\} \\
 & = & \left(\frac{p}{p_{Te}}\right)^{2n}\xi_{0}^{2}\left(1-\xi_{0}^{2}\right)^{(n-1)}\frac{v_{Te}}{v|\xi_{0}|}\\
 &  & \times\left\{ \frac{\xi}{\xi_{0}}\Psi^{n}\exp\left(-\frac{\left(k_{\Vert,\text{res}}-k_{\Vert,0}\right)^{2}}{\Delta k_{\Vert}^{2}}\right)2\pi\delta(\theta-\theta_{b})D_{0}^{\text{EC-O}(n)}\right\} \\
 & = & \left(\frac{p}{p_{Te}}\right)^{2n}\xi_{0}^{2}\left(1-\xi_{0}^{2}\right)^{(n-1)}\frac{v_{Te}}{v|\xi_{0}|}\\
 &  & \frac{1}{\lambda}\left[\frac{1}{2}\sum_{\sigma=\pm1}\right]_{T}\int_{-\theta_{c}}^{\theta_{c}}\frac{d\theta}{2\pi}\Psi^{n}\exp\left(-\frac{\left(k_{\Vert,\text{res}}-k_{\Vert,0}\right)^{2}}{\Delta k_{\Vert}^{2}}\right)2\pi\delta(\theta-\theta_{b})D_{0}^{\text{EC-O}(n)}\end{eqnarray*}
 et finalement \begin{eqnarray*}
D_{(0)}^{\text{EC-O}(n)}(p,\xi_{0}) & \equiv & \left\{ \Psi D^{\text{EC-O}(n)}(p,\xi)\right\} \\
 & = & \left(\frac{p}{p_{Te}}\right)^{2n}\xi_{0}^{2}\left(1-\xi_{0}^{2}\right)^{(n-1)}\frac{v_{Te}}{v|\xi_{0}|}\\
\times\frac{1}{\lambda}\left[\frac{1}{2}\sum_{\sigma=\pm1}\right]_{T}H(|\theta_{b}| & < & \theta_{c})\Psi_{b}^{n}\exp\left(-\frac{\left(k_{\Vert,\text{res},b}-k_{\Vert,0}\right)^{2}}{\Delta k_{\Vert}^{2}}\right)D_{0,b}^{\text{EC-O}(n)}\end{eqnarray*}
 avec \[
D_{0,b}^{\text{EC-O}(n)}=\lim_{V\rightarrow\infty}\frac{\pi q_{e}^{2}\Delta\theta_{b}}{V}\int\frac{d^{2}k_{\bot}}{(2\pi)^{3}}\frac{1}{n!^{2}}\left(N_{\bot,b}\frac{p_{Te}}{m_{e}c}\frac{\omega_{b}}{2\omega_{ce}}\right)^{2n}\left|E_{\mathbf{k},\Vert,0,b}\right|^{2}(\overline{k}_{\bot})\frac{1}{v_{Te}}\]

